\chapter{تحلیل ریسک امنیتی}
\label{chap:security}
%=====================================================================

\section{آسیب‌پذیری‌های شناخته‌شده}

\begin{enumerate}
\item \textbf{آسیب‌پذیری گرهٔ مورد اعتماد:} همان‌طور که در بخش \ref{sec:trustednode} و \ref{sec:bankingtrustednode} بحث شد، هر گرهٔ مورد اعتماد به‌طور ذاتی یک نقطهٔ شکست امنیتی است؛ حتی معماری چین با $32$ گره نیز از این آسیب‌پذیری ساختاری مبرا نیست \Cite{iiss2026governance}.
\item \textbf{حملات کانال جانبی سخت‌افزاری:} حملاتی نظیر «کوری‌سازی آشکارساز» \lr{(Detector-Blinding Attack)} با نور شدید، آشکارساز تک‌فوتونی را به رفتار کلاسیک وادار کرده و امکان استراق‌سمع بدون افزایش قابل‌کشف \lr{QBER} را فراهم می‌کند. راهکار: استفاده از \lr{MDI-QKD} برای حذف اعتماد به آشکارساز.
\item \textbf{عدم‌قطعیت اثبات امنیتی دستگاه‌های واقعی:} فرمول \ref{eq:gllp} فرض می‌کند دستگاه‌ها دقیقاً مطابق مدل نظری عمل می‌کنند؛ هرگونه انحراف سخت‌افزاری واقعی می‌تواند شکاف امنیتی ایجاد کند -- دلیل اصلی توصیهٔ ممیزی مستقل و آزمون نفوذ منظم.
\item \textbf{فقدان استاندارد جهانی یکپارچهٔ گواهی‌دهی:} برخلاف \lr{PQC} که استانداردهای \lr{NIST} را دارد، \lr{QKD} هنوز از پراکندگی نهادی در گواهی‌دهی رنج می‌برد \Cite{iiss2026governance}، که ریسک تدارکاتی (انتخاب فروشنده) را افزایش می‌دهد.
\item \textbf{ریسک الگوریتمی در لایهٔ \lr{PQC}:} اگرچه \lr{ML-KEM} و \lr{ML-DSA} اکنون استاندارد نهایی محسوب می‌شوند، سابقهٔ شکست \lr{SIKE} و \lr{Rainbow} در همان فرآیند ارزیابی \lr{NIST} \Cite{nistir8547} نشان می‌دهد که هیچ فرض ریاضی‌ای کاملاً مصون از تجدیدنظر نیست -- دلیل اصلی توصیهٔ راهبرد ترکیبی فصل \ref{chap:hybrid}.
\end{enumerate}

\section{راهکارهای کاهش ریسک پیشنهادی}

\begin{itemize}
\item الزام معماری \textbf{ترکیبی \lr{QKD+PQC}} برای تمام کاربردهای حساس بانکی از روز نخست، مطابق فصل \ref{chap:hybrid}؛
\item استقرار \lr{MDI-QKD} برای گره‌های با بالاترین حساسیت (بانک مرکزی و مرکز تسویهٔ بین‌بانکی)؛
\item تدوین سند داخلی «الزامات امنیت فیزیکی گرهٔ مورد اعتماد» با همکاری مرکز افتای ریاست‌جمهوری و بانک مرکزی، پیش از عملیاتی‌سازی؛
\item برگزاری آزمون نفوذ سالانه و ممیزی مستقل بین‌المللی (در حد امکان با رعایت الزامات حاکمیتی) برای اعتبارسنجی سطح امنیت واقعی در برابر مدل نظری؛
\item طراحی معماری با \textbf{چابکی رمزنگاری \lr{(Crypto-Agility)}} به‌گونه‌ای که جایگزینی الگوریتم \lr{PQC} در آینده (در صورت کشف ضعف) بدون بازطراحی کامل زیرساخت امکان‌پذیر باشد.
\end{itemize}

%=====================================================================
